\documentclass{article}
\usepackage{../style/dhruv_theory_style}

% -------------------------------------------------------
% Document metadata
% -------------------------------------------------------
\title{\textbf{Part B: Theory Problem}\\[4pt]
       \large $\Delta$-Separated Threshold-Realizable Sequences}
\author{DSC 180}
\date{\today}

% =======================================================
\begin{document}
% =======================================================

\maketitle
\tableofcontents
\newpage

% =======================================================
\section{Problem Overview}
% =======================================================

This problem starts from the threshold story in lecture and then pushes it one
step further.

Throughout, let $\mathcal{X} = [0,1]$, $\mathcal{Y} = \{-1,+1\}$, and let
\[
h_\theta(x) = \sign(x-\theta),
\]
where we use the convention $\sign(z)=+1$ for $z \ge 0$ and $\sign(z)=-1$ for
$z<0$.

We will study sequences with one extra structural assumption beyond what was
explicitly analyzed in lecture.
\newline

\textbf{$\Delta$-Separated Threshold-Realizable Sequences.}
We say a sequence $((x_t,y_t))_{t=1}^T$ is \emph{$\Delta$-separated
threshold-realizable} if there exist $\theta^* \in [0,1]$ and $\Delta > 0$
such that
\[
y_t = h_{\theta^*}(x_t)
\qquad \text{and} \qquad
|x_t-\theta^*| \ge \Delta
\qquad \text{for every round } t.
\]

\qquad
\qquad

% =======================================================
\section{From Continuous Thresholds to a Finite Class}
% =======================================================

\textbf{Question.}
Design a finite grid $G \subseteq [0,1]$ whose size depends only on $\Delta$,
and consider the associated finite threshold class
\[
H_G = \{h_\theta : \theta \in G\}.
\]
Prove that for every $\Delta$-separated threshold-realizable sequence, there
exists some $\theta \in G$ such that
\[
h_\theta(x_t)=y_t
\qquad \text{for all } t=1,\dots,T.
\]
Then use the Halving theorem from lecture to derive a mistake bound of order
\[
O\!\left(\log\left(\frac{1}{\Delta}\right)\right)
\]
for an explicit online learner.

Your answer should clearly state:
\begin{itemize}[leftmargin=2.5em]
    \item your choice of grid $G$,
    \item the size of $G$ as a function of $\Delta$,
    \item the resulting mistake bound.
\end{itemize}

\begin{solutionbox}

\subsection{Construction of the Grid}

We choose a grid with spacing $\Delta/2$:
\[
G
=
\left\{
k\cdot \frac{\Delta}{2}
:
k=0,1,2,\dots,\left\lfloor \frac{2}{\Delta}\right\rfloor
\right\}
\cap [0,1].
\]

Its size satisfies
\[
|G|
\le
\left\lfloor \frac{2}{\Delta}\right\rfloor + 1
=
O\!\left(\frac{1}{\Delta}\right).
\]

The reason for spacing $\Delta/2$ is that for every $\theta^* \in [0,1]$,
there exists some grid point $\theta \in G$ such that
\[
|\theta-\theta^*| \le \frac{\Delta}{2}.
\]

\subsection{Proof That a Grid Threshold Is Consistent}

\begin{claim}
\leavevmode\\
    For every $\Delta$-separated threshold-realizable sequence
    $(x_t,y_t)_{t=1}^T$, there exists $\theta \in G$ such that
    \[
    h_\theta(x_t)=y_t
    \qquad \text{for all } t=1,\dots,T.
    \]
\end{claim}
\begin{proof}
\leavevmode\\
Because the sequence is $\Delta$-separated threshold-realizable, there exists
$\theta^* \in [0,1]$ such that
\[
y_t = h_{\theta^*}(x_t)=\sign(x_t-\theta^*)
\qquad \text{for all } t,
\]
and
\[
|x_t-\theta^*| \ge \Delta
\qquad \text{for all } t.
\]

By construction of the grid, there exists some $\theta \in G$ with
\[
|\theta-\theta^*| \le \frac{\Delta}{2}.
\]

We now show that $h_\theta(x_t)=h_{\theta^*}(x_t)$ for every $t$.

\textbf{Case 1:} $y_t=+1$.

Then $x_t \ge \theta^*$, and by separation,
\[
x_t-\theta^* \ge \Delta.
\]
Using $|\theta-\theta^*| \le \Delta/2$,
\[
x_t-\theta
=
(x_t-\theta^*)-(\theta-\theta^*)
\ge
\Delta-\frac{\Delta}{2}
=
\frac{\Delta}{2}
>0.
\]
So $x_t \ge \theta$, hence
\[
h_\theta(x_t)=+1=y_t.
\]

\textbf{Case 2:} $y_t=-1$.

Then $x_t<\theta^*$, and by separation,
\[
\theta^*-x_t \ge \Delta.
\]
Again using $|\theta-\theta^*| \le \Delta/2$,
\[
\theta-x_t
=
(\theta^*-x_t)-(\theta^*-\theta)
\ge
\Delta-\frac{\Delta}{2}
=
\frac{\Delta}{2}
>0.
\]
So $x_t<\theta$, hence
\[
h_\theta(x_t)=-1=y_t.
\]

Thus $h_\theta(x_t)=y_t$ for every $t$.
\end{proof}

\subsection{Halving Algorithm and Mistake Bound}

Consider the finite class
\[
H_G=\{h_\theta:\theta\in G\}.
\]
Run the Halving algorithm on this class.

Since the claim above shows that there exists a hypothesis in $H_G$ consistent
with the entire sequence, the Halving theorem applies. Therefore the number of
mistakes is at most
\[
M \le \log_2 |H_G| = \log_2 |G|.
\]
Using the bound on the size of the grid,
\[
M
\le
\log_2\!\left(\left\lfloor \frac{2}{\Delta}\right\rfloor+1\right)
=
O\!\left(\log\left(\frac{1}{\Delta}\right)\right).
\]

\subsection*{Final Answer for Part 1}

\begin{itemize}[leftmargin=2.5em]
    \item \textbf{Choice of grid:}
    \[
    G=
    \left\{
    k\cdot \frac{\Delta}{2}
    :
    k=0,1,2,\dots,\left\lfloor \frac{2}{\Delta}\right\rfloor
    \right\}
    \cap [0,1].
    \]

    \item \textbf{Size of the grid:}
    \[
    |G|
    \le
    \left\lfloor \frac{2}{\Delta}\right\rfloor+1
    =
    O\!\left(\frac{1}{\Delta}\right).
    \]

    \item \textbf{Mistake bound:}
    \[
    M
    \le
    \log_2\!\left(\left\lfloor \frac{2}{\Delta}\right\rfloor+1\right)
    =
    O\!\left(\log\left(\frac{1}{\Delta}\right)\right).
    \]
\end{itemize}

\end{solutionbox}

% =======================================================
\section{A Positive Margin from Separation}
% =======================================================

\textbf{Question.}
\newline
View thresholds as linear predictors by choosing a feature map
\[
\phi:[0,1]\to \mathbb{R}^d
\]
and a unit vector $u^*$ depending on $\theta^*$.

Prove that every $\Delta$-separated threshold-realizable sequence is linearly
separable with margin at least $c\Delta$ for some absolute constant $c>0$
under your representation, while
\[
\|\phi(x)\|\le R
\qquad \text{for all } x\in [0,1]
\]
for some absolute constant $R$.

Then use the Perceptron theorem from lecture to derive an explicit mistake
bound of order
\[
O\!\left(\frac{1}{\Delta^2}\right).
\]

Your answer should clearly specify your chosen feature map, the margin lower
bound, the norm bound, and the final mistake bound.

\begin{solutionbox}

\subsection{Feature Map and Linear Representation}

Chosen
\[
\phi(x)=(x,1)\in \R^2.
\]
Let
\[
w^*=(1,-\theta^*).
\]
Then
\[
\langle w^*,\phi(x)\rangle
=
1\cdot x + (-\theta^*)\cdot 1
=
x-\theta^*.
\]
Therefore
\[
\sign(\langle w^*,\phi(x)\rangle)
=
\sign(x-\theta^*)
=
h_{\theta^*}(x).
\]

So this representation realizes threshold classification exactly.

Now normalize:
\[
u^*=\frac{w^*}{\|w^*\|}.
\]
Since $\theta^*\in[0,1]$,
\[
\|w^*\|
=
\sqrt{1+(\theta^*)^2}
\le
\sqrt{2}.
\]

\subsection{Margin Lower Bound}

For every round $t$,
\[
y_t\langle u^*,\phi(x_t)\rangle
=
\frac{y_t\langle w^*,\phi(x_t)\rangle}{\|w^*\|}
=
\frac{y_t(x_t-\theta^*)}{\|w^*\|}.
\]
Because
\[
y_t=\sign(x_t-\theta^*),
\]
we get
\[
y_t(x_t-\theta^*) = |x_t-\theta^*|.
\]
By the $\Delta$-separation assumption,
\[
|x_t-\theta^*| \ge \Delta.
\]
Also $\|w^*\|\le \sqrt{2}$. Hence
\[
y_t\langle u^*,\phi(x_t)\rangle
\ge
\frac{\Delta}{\sqrt{2}}.
\]

Thus the sequence is linearly separable with margin at least
\[
\gamma \ge \frac{\Delta}{\sqrt{2}}.
\]
So we can take
\[
c=\frac{1}{\sqrt{2}}.
\]

\subsection{Norm Bound}

For any $x\in[0,1]$,
\[
\|\phi(x)\|^2 = x^2+1 \le 2.
\]
Therefore
\[
\|\phi(x)\| \le \sqrt{2}.
\]
So we may take
\[
R=\sqrt{2}.
\]

\subsection{Perceptron Mistake Bound}

By the Perceptron theorem, the number of mistakes satisfies
\[
M \le \left(\frac{R}{\gamma}\right)^2.
\]
Substituting the bounds above,
\[
M
\le
\left(\frac{\sqrt{2}}{\Delta/\sqrt{2}}\right)^2
=
\left(\frac{2}{\Delta}\right)^2
=
\frac{4}{\Delta^2}.
\]
Hence
\[
M = O\!\left(\frac{1}{\Delta^2}\right).
\]

\subsection*{Final Answer for Part 2}

\begin{itemize}[leftmargin=2.5em]
    \item \textbf{Feature map:}
    \[
    \phi(x)=(x,1)\in\R^2.
    \]

    \item \textbf{Margin lower bound:}
    \[
    \gamma \ge \frac{\Delta}{\sqrt{2}}.
    \]

    \item \textbf{Norm bound:}
    \[
    \|\phi(x)\|\le \sqrt{2}
    \qquad \text{for all } x\in[0,1].
    \]

    \item \textbf{Mistake bound:}
    \[
    M\le \frac{4}{\Delta^2}
    =
    O\!\left(\frac{1}{\Delta^2}\right).
    \]
\end{itemize}

\end{solutionbox}

% =======================================================
\section{Comparison and Interpretation}
% =======================================================

\textbf{Question.}
\newline
Explain why the continuous-threshold impossibility phenomenon from lecture does
not contradict Part 1.

Then compare the two mistake bounds you obtained in Parts 1 and 2. Why do they
scale differently with $\Delta$? What is each argument measuring about the
problem?

\begin{solutionbox}

The continuous-threshold impossibility result from lecture applies to the full
class
\[
\mathcal{H}=\{h_\theta:\theta\in[0,1]\}
\]
with no structural restriction on the sequence. In that unrestricted setting,
an adversary can place examples arbitrarily close to the true threshold, which
prevents an online learner from obtaining a finite worst-case mistake bound.

Part 1 does not contradict this because it assumes something extra:
$\Delta$-separation. Once every example must remain at least $\Delta$ away from
the true threshold $\theta^*$, thresholds within distance $\Delta/2$ of
$\theta^*$ become indistinguishable on the observed data. That is why the
continuous class can be replaced by a finite grid. So the impossibility result
and the positive result in Part 1 refer to different settings.

Now compare the two bounds:
\[
\text{Part 1: } O\!\left(\log\left(\frac{1}{\Delta}\right)\right),
\qquad
\text{Part 2: } O\!\left(\frac{1}{\Delta^2}\right).
\]

They scale differently because they come from different kinds of arguments.

The Part 1 argument is \emph{combinatorial}. It counts how many effectively
different hypotheses remain after using the separation assumption. Since the
effective finite class has size
\[
O\!\left(\frac{1}{\Delta}\right),
\]
the Halving algorithm gives a logarithmic mistake bound.

The Part 2 argument is \emph{geometric}. It uses linear separability and the
margin bound. The Perceptron theorem depends on
\[
\left(\frac{R}{\gamma}\right)^2.
\]
Because the margin is proportional to $\Delta$, the mistake bound scales like
$1/\Delta^2$.

So the Halving argument measures how many candidate thresholds need to be
distinguished, while the Perceptron argument measures how hard the problem is
geometrically in terms of margin.

\end{solutionbox}

% =======================================================
\section{Optional Strengthening: Audit an AI Proof}
% =======================================================

\textbf{Question.}
An AI assistant gives the following argument:

\hangindent=2em
\hangafter=0
Because every example is at least $\Delta$ away from the true threshold, continuous thresholds effectively form a class of size $O(1/\Delta)$. Therefore any online learner, including Perceptron, must make at most $O(\log(1/\Delta))$ mistakes on every $\Delta$-separated threshold-realizable sequence.

Identify at least two mathematical problems with this argument. Then write a
corrected statement that is true and prove it.

\begin{solutionbox}

There are at least two mathematical problems with the AI's argument.

\textbf{Problem 1.}
It incorrectly concludes that because the effective class size is
$O(1/\Delta)$, \emph{any} online learner must make only
$O(\log(1/\Delta))$ mistakes. That is not true. A class-size argument gives a
mistake bound only for specific algorithms such as the Halving algorithm. It is
not a universal guarantee for every possible learner.

\textbf{Problem 2.}
It incorrectly applies the same conclusion to Perceptron. The Perceptron
theorem does not use class size. It uses margin and feature norm. In this
problem, the Perceptron bound is
\[
O\!\left(\frac{1}{\Delta^2}\right),
\]
not
\[
O\!\left(\log\left(\frac{1}{\Delta}\right)\right).
\]

\textbf{Corrected statement.}
For every $\Delta>0$, if we run the Halving algorithm on the finite threshold
class
\[
H_G=\{h_\theta:\theta\in G\},
\]
where
\[
G=
\left\{
k\cdot \frac{\Delta}{2}
:
k=0,1,2,\dots,\left\lfloor \frac{2}{\Delta}\right\rfloor
\right\}
\cap [0,1],
\]
then on every $\Delta$-separated threshold-realizable sequence, the number of
mistakes is at most
\[
\log_2\!\left(\left\lfloor \frac{2}{\Delta}\right\rfloor+1\right)
=
O\!\left(\log\left(\frac{1}{\Delta}\right)\right).
\]

\textbf{Proof.}
From Part 1, there exists some $\theta\in G$ such that
\[
h_\theta(x_t)=y_t
\qquad \text{for all } t=1,\dots,T.
\]
So the class $H_G$ contains a hypothesis consistent with the entire sequence.

The Halving algorithm starts with all hypotheses in $H_G$. Whenever it makes a
mistake, at least half of the currently surviving hypotheses are eliminated.
Because a consistent hypothesis is never removed, the version space always
contains at least one hypothesis. Therefore after $M$ mistakes,
\[
1 \le \frac{|G|}{2^M}.
\]
This implies
\[
M \le \log_2 |G|.
\]
Using
\[
|G| \le \left\lfloor \frac{2}{\Delta}\right\rfloor+1,
\]
we obtain
\[
M
\le
\log_2\!\left(\left\lfloor \frac{2}{\Delta}\right\rfloor+1\right)
=
O\!\left(\log\left(\frac{1}{\Delta}\right)\right).
\]

This corrected statement is true for the Halving algorithm specifically. It
does not imply the same bound for arbitrary online learners or for Perceptron.

\end{solutionbox}

\end{document}